\section{A smooth integral cubic and its Tate multiplicity}
\label{sec:cubic}
This section places the projective coherent-line construction beside an affine construction over $\mathbb Z$ in which the exact count is a multiplicity in a motive, a Chow group, and a Grothendieck group. The geometric suspension mechanism is due to Asok, Dubouloz, and \O stv\ae r (2023, Theorem~3.1.1 and Remark~3.1.3; preprint pp.~8--9), whose Construction~4.2.1 already includes this form of equation \cite{ff-asok}. The arithmetic encoding and its counting interpretation are given explicitly below.

This construction has a different output from two established universality results. Lee and Vakil (2013, Theorem~1.3, preprint p.~2) show that the disjoint union of their point--line incidence schemes realizes every finite-type singularity over $\mathbb Z$, up to their equivalence generated by smooth morphisms \cite{ff-lee-vakil}. Their schemes impose both incidences and nonincidences, together with distinctness and intersection conditions. The theorem does not prescribe a count of finite-field points. Grochow and Qiao (2023, Theorem~A, Corollary~A, and Theorem~B) give polynomial reductions among tensor, algebra, group, and form isomorphism problems, including a reduction from fixed-order tensor isomorphism to three-tensor isomorphism; their cubic-form equivalence statement assumes characteristic different from two and three \cite{ff-grochow-qiao}. That theorem preserves orbit equivalence. Here the preserved datum is the number of accepting words, which appears as a specified cohomological multiplicity on a smooth geometrically integral scheme.

Start with a Boolean verification circuit. Introduce a variable for each primary input and each gate wire, so that the total number is $v\geq1$. Include every Boolean equation $x_j^2-x_j=0$. Encode an AND gate by $z-xy=0$, a NOT gate by $z+x-1=0$, and the acceptance condition by $o-1=0$. These gates suffice after replacing other Boolean operations by constant-size circuits. Let the full list be $f_1,\ldots,f_h$, with $f_1=x_1^2-x_1$. Each equation has degree at most two, at most three terms, and integral coefficients of bounded size.

The Boolean quotient is
\[
 \mathbb Z[x_1,\ldots,x_v]/(x_j^2-x_j:1\leq j\leq v)
      \simeq \prod_{a\in\{0,1\}^v}\mathbb Z.
\]
On every factor each gate equation evaluates to $0$, $1$, or $-1$. Quotienting deletes the illegal factors and keeps the legal ones. Each accepting primary word has exactly one extension to its wires, so
\[
 Y:=V(f_1,\ldots,f_h)\simeq\coprod_{S}\operatorname{Spec}\mathbb Z.
\]
This is a scheme-theoretic identity over the integer base, stronger than equality of rational-point sets over one field.

For a finite-type scheme $T$ over $\mathbb Z$, write $[T]$ for its class in the Grothendieck ring of such schemes, defined by the scissor relation $[T]=[W]+[T\setminus W]$ for a closed subscheme $W$, with multiplication induced by fibre product over $\mathbb Z$.

\begin{theorem}\label{thm:cubic}
The affine hypersurface
\begin{equation}\label{eq:cubic}
 X_\Phi:\quad t^2-t=\sum_{i=1}^h y_i f_i(x)
 \quad\subset\mathbb A^{v+h+1}_{\mathbb Z}
\end{equation}
is smooth over $\mathbb Z$, with geometrically integral fibers of dimension $v+h$, and has a sparse cubic equation of polynomial bit size. Writing $\mathbb L=[\mathbb A^1_{\mathbb Z}]$, one has
\[
 [X_\Phi]=\mathbb L^{v+h}+S\mathbb L^h,
 \qquad
 \#X_\Phi(\mathbb F_q)=q^{v+h}+Sq^h
\]
for every prime power $q$. For every field $k$ and $\lambda\neq\operatorname{char}k$, its geometric compactly supported cohomology is
\[
 H_c^{2h}(X_{\Phi,\bar k},\mathbb Q_\lambda)
        =\mathbb Q_\lambda(-h)^S,
 \qquad
 H_c^{2(v+h)}(X_{\Phi,\bar k},\mathbb Q_\lambda)
        =\mathbb Q_\lambda(-v-h),
\]
with all other groups zero, compatibly with the Galois action. Consequently
\[
 Z(X_\Phi/\mathbb F_q,T)
      =\frac{1}{(1-q^{v+h}T)(1-q^hT)^S},
 \qquad
 \zeta_{X_\Phi}(s)=\zeta(s-v-h)\zeta(s-h)^S.
\]
\end{theorem}
\begin{proof}
Write $G=t^2-t-\sum y_if_i$. In the coordinate ring,
\[
 (2t-1)\frac{\partial G}{\partial t}
        +4\sum_i y_i\frac{\partial G}{\partial y_i}
 =(2t-1)^2-4\sum_i y_if_i=1.
\]
The relative partial derivatives in the $t,y$ coordinates generate the unit ideal. Thus projection $\pi:X_\Phi\to\mathbb A^v_{\mathbb Z}$ is smooth of relative dimension $h$, including in characteristic two. Composition with the smooth affine base proves the first smoothness assertion.

Over an algebraic closure $\bar k$, view $G$ as a polynomial in $y_1$ over $R=\bar k[x,t,y_2,\ldots,y_h]$. Its leading coefficient is the nonzero polynomial $-f_1(x)$, and its constant coefficient is monic in $t$. They have no common irreducible factor: a factor of $f_1$ depends only on $x$ and cannot divide a polynomial with leading $t$-coefficient one. Thus $G$ is primitive over the UFD $R$. It is irreducible of degree one in $y_1$ over $\operatorname{Frac}(R)$, so Gauss's lemma makes the quotient an integral domain. This proves geometric integrality.

Stratify $\mathbb A^v$ by $Y$ and the locally closed loci where $f_i$ is the first nonzero member of the list. Over the latter locus solve uniquely for $y_i$; the other $h-1$ variables $y_j$ and $t$ are free. The fiber is therefore $\mathbb A^h$. Over $Y$ the equation is $t(t-1)=0$, giving two disjoint copies of $Y\times\mathbb A^h$. Since $[Y]=S$, additivity yields the class and point-count identities.

The cohomology statement requires more than a class calculation. By compact-support base change, the geometric stalks of $R^j\pi_!\mathbb Q_\lambda$ are zero except for $j=2h$, because every fiber is either one or two affine $h$-spaces. Write $\mathcal T=R^{2h}\pi_!\mathbb Q_\lambda$. The involution $\iota:t\mapsto1-t$ yields the projectors $(1\pm\iota)/2$ over $\mathbb Q_\lambda$, also when $\lambda=2$. Off $Y$, it acts trivially on the top compact cohomology of the affine-space fiber. Over $Y$ it exchanges the two labeled components.

The relative trace $\mathcal T(h)\to\mathbb Q_\lambda$ restricts to an isomorphism on the plus eigensheaf: on a single component it is the affine-space trace, and on the diagonal in two components it is multiplication by two. The minus eigensheaf has zero stalks off $Y$ and the anti-diagonal line over every point of $Y$. With $i:Y\hookrightarrow\mathbb A^v$, localization therefore identifies it with $i_*\mathbb Q_\lambda(-h)$. These stalkwise identifications are isomorphisms of sheaves and are defined over the ground field. Hence
\[
 R\pi_!\mathbb Q_\lambda
   \simeq(\mathbb Q_\lambda\oplus i_*\mathbb Q_\lambda)(-h)[-2h].
\]
Taking compact supports on the base gives the two displayed groups. The ingredients are compact-support base change, the relative trace, and affine-space cohomology in SGA~4, Expos\'e~XVII, Theorems~5.2.6 and~5.4.3, and Expos\'e~XVIII, (2.8.1) and Theorem~2.9 \cite{ff-sga4}. No cancellation in a virtual representation is being used. The zeta identities follow from the all-extension-field counts; the arithmetic Euler product converges initially for $\operatorname{Re}s>v+h+1$.
\end{proof}

The construction gives a numerical one-call reduction from exact witness counting to exact point counting over each fixed finite field:
\[
 S=\frac{\#X_\Phi(\mathbb F_q)-q^{v+h}}{q^h}.
\]
The numerator is known to be divisible by $q^h$, and all integers involved have polynomially many bits. For a modulus $m$ coprime to $q$, the same equation recovers $S\bmod m$ by modular inversion. For $m=p^r$ and $q=p$, the first $h$ digits of the count vanish; precision $p^{h+r}$ recovers $S\bmod p^r$ after subtraction and exact division. Thus the rank-one coherent theorem and the cubic theorem locate the hard count in different arithmetic places.

\subsection{Motivic spheres and affine-space recognition}
For a smooth scheme $T$ over a perfect field, $M(T)$ denotes its effective geometric motive with integral coefficients. The notation $\mathbb Z(v)[2v]$ denotes the $v$th Tate twist of the unit motive, shifted by $2v$. We write $\operatorname{CH}^j(T,r)$ for Bloch's higher Chow group of codimension $j$ and homological degree $r\geq0$. When $a\in Y$ indexes a wedge or a generator below, it ranges over the $S$ accepting Boolean words, equivalently the $S$ components $\operatorname{Spec}\mathbb Z$ of $Y$.

\begin{theorem}\label{thm:motivic}
Point $X_\Phi$ at $x=t=y=0$. In the pointed unstable $\mathbb A^1$-homotopy category over $\mathbb Z$,
\[
 (X_\Phi,0)\simeq\bigvee_{a\in Y}(\mathbb P^1,\infty)^{\wedge v}.
\]
The empty wedge is the pointed base. The equivalence holds after every base change with the same indexing and framing. Over a perfect field $k$ it gives
\[
 M(X_{\Phi,k})\simeq\mathbb Z\oplus\mathbb Z(v)[2v]^S,
\]
and
\[
 \operatorname{CH}^j(X_{\Phi,k},r)
 \simeq\operatorname{CH}^j(k,r)
       \oplus\operatorname{CH}^{j-v}(k,r)^S.
\]
Over $\mathbb C$, $X_\Phi(\mathbb C)$ has homotopy type $\bigvee_S S^{2v}$. Moreover $S=0$ if and only if $X_\Phi\simeq\mathbb A^{v+h}_{\mathbb Z}$, and if and only if its fiber over any specified field is isomorphic to affine space.
\end{theorem}
\begin{proof}
The smooth projection $\pi$ has section $s(x)=(x,0,0)$. On $U=\mathbb A^v\setminus Y$ it is locally trivial with fiber $\mathbb A^h$, by solving for a $y_i$. Hence $\pi^{-1}(U)\to U$ is an $\mathbb A^1$-weak equivalence. The closed scheme $Y$ is smooth over $\mathbb Z$. Its immersion in $\mathbb A^v$ has trivial normal bundle of rank $v$: the Boolean equations frame the conormal bundle because on the component indexed by $a$,
\[
 d(x_j^2-x_j)=(2a_j-1)\,dx_j,
\]
and each multiplier is an integral unit. Smooth base change pulls this frame back to $X_Y\subset X$.

The homotopy-purity cofiber theorem of Asok, Dubouloz, and \O stv\ae r, Theorem~3.1.1, applies to this smooth map, its section, and the contractible base. It gives
\[
 \operatorname{Th}(N_{Y/\mathbb A^v})\longrightarrow
 \operatorname{Th}(N_{X_Y/X})\longrightarrow(X,s(0)).
\]
Now $X_Y=D_0\amalg D_1$ with $D_e=Y\times\mathbb A^h$ given by $t=e$, and the section lands in $D_0$. Contracting the affine factors and using the compatible normal frames identifies the first map with inclusion of the first summand
\[
 \bigvee_S(\mathbb P^1)^{\wedge v}\longrightarrow
 \left(\bigvee_S(\mathbb P^1)^{\wedge v}\right)
 \vee\left(\bigvee_S(\mathbb P^1)^{\wedge v}\right).
\]
Its cofiber is the second summand, proving the equivalence. This explicitly checks the map of Thom spaces, which is indispensable for the wedge count. Over a perfect field, the reduced motive of the pointed $v$-fold projective-line smash is $\mathbb Z(v)[2v]$; the rational point splits the unit motive. Motivic cohomology and its higher-Chow comparison give the stated groups \cite[Theorem~19.1, p.~159]{ff-mvw}. Complex realization sends this motivic sphere to $S^{2v}$, giving the topological assertion.

If $S=0$, the ideal $(f_1,\ldots,f_h)$ is the whole integral polynomial ring, since its quotient is the empty product, namely the zero ring. Choose a column $a(x)$ with $\sum a_if_i=1$. The maps
\[
 z=y+a(2t-t^2),\qquad
 t=\sum_i f_i z_i,\qquad y=z+a(t^2-2t)
\]
are inverse polynomial maps between $X$ and $\mathbb A^{v+h}_{\mathbb Z}$. Indeed the first satisfies $f\cdot z=t$, and the second satisfies $f\cdot y=t^2-t$. No freeness assertion for an arbitrary projective module is needed. If $S>0$, the nonzero reduced motivic or geometric cohomology forbids an affine-space isomorphism over any field; base change to its algebraic closure supplies the obstruction. The same argument rules out the integral isomorphism. Finally the complex sphere wedge is contractible exactly for $S=0$.
\end{proof}

There is a useful effective distinction in this last assertion. If the verifier is unsatisfiable, the multipliers $a_i$ can be chosen of degree at most $v$ and coefficient bit length $O(v)$. For each Boolean word $b$ choose an equation taking value $\epsilon_b\in\{1,-1\}$ and set
\[
 e_b(x)=\prod_{b_j=1}x_j\prod_{b_j=0}(1-x_j),\qquad
 a_i^{(0)}=\sum_{b:\,i(b)=i}\epsilon_b e_b.
\]
The polynomial $\sum a_i^{(0)}f_i-1$ has zero multilinear remainder modulo the Boolean equations. Telescoping each monomial reduces it integrally by $x_j^2-x_j$ with quotient degree at most $v$. Subtracting those quotients gives $a_i$. Since $\sum_b\|e_b\|_1=3^v$ and each $f_i$ has coefficient norm at most three, all coefficients have $O(v)$ bits. Nevertheless there may be exponentially many monomials. These degree and height bounds do not give a polynomial-time isomorphism finder or a short certificate of unsatisfiability.

The affine-space criterion in Theorem~\ref{thm:motivic} depends on this particular family. Katz and Rodriguez Villegas (2026, Example~2.3, pp.~4--5) discuss the smooth Russell threefold $x^2y+x+z^2+t^3=0$, which has exactly $q^3$ points over every finite field but is not isomorphic to $\mathbb A^3$ over $\mathbb C$ \cite{ff-katz-rv}. Their Section~3 also gives a smooth irreducible polynomial-count variety with nondiagonal Hodge contributions that cancel between cohomological degrees. The decomposition proved above identifies the cohomology groups themselves, and the explicit polynomial inverse proves affine-space recognition in the zero-count case; neither conclusion is inferred solely from the counting polynomial.

\subsection{The Chow ring and the integral gamma filtration}
The Chow ring can be calculated over every field without appealing to the perfect-field motive. For a satisfying wire word $a$ and $e\in\{0,1\}$, set
\[
 D_{a,e}=V(x-a,t-e)\subset X,\qquad D_{a,e}\simeq\mathbb A^h.
\]
Let $u_a$ be the codimension-$v$ cycle class $[D_{a,0}]$.

In the $K$-theory statements, $K_0(X)$ is the Grothendieck ring of algebraic vector bundles and $I$ is the kernel of its rank map to $\mathbb Z$. Write $\psi^m$ for the $m$th Adams operation, $m\geq1$, and
\[
 \lambda_t(z)=\sum_{j\geq0}\lambda^j(z)t^j,\qquad
 \gamma_t(z)=\lambda_{t/(1-t)}(z)=\sum_{j\geq0}\gamma^j(z)t^j
\]
for the total exterior-power and gamma operations in the formal variable $t$. The gamma filtration $F_\gamma^iK_0$ is generated by products of $\gamma^{j_\nu}(z_\nu)$ with $z_\nu\in I$ and $\sum_\nu j_\nu\geq i$; its associated graded is $\operatorname{gr}_\gamma^iK_0=F_\gamma^iK_0/F_\gamma^{i+1}K_0$. The braces $\left\{\begin{smallmatrix}v\\j\end{smallmatrix}\right\}$ denote Stirling numbers of the second kind.

\begin{proposition}\label{prop:chow-gamma}
For every field $k$,
\[
 \operatorname{CH}^*(X_{\Phi,k})
  \simeq\mathbb Z[u_a:a\in Y]/(u_au_b:a,b\in Y),
 \qquad\deg u_a=v,
\]
and $[D_{a,1}]=-u_a$. The cycle map $\operatorname{CH}^j(X)\otimes\mathbb Q_\lambda\to H^{2j}(X_{\bar k},\mathbb Q_\lambda(j))$ is an isomorphism for every $j$; odd cohomology is zero. Over $\mathbb Z$ and after base change to every field,
\[
 K_0(X)=\mathbb Z\oplus I,\qquad I\simeq\mathbb Z^S,\qquad I^2=0,
 \qquad\psi^m(z)=m^v z\quad(z\in I).
\]
In this augmented $\lambda$-ring,
\begin{align}
 \lambda_t(z)&=1+z\sum_{m\geq1}(-1)^{m-1}m^{v-1}t^m,\\
 \gamma_t(z)&=1+z\sum_{j=1}^v(-1)^{j-1}(j-1)!
                          \left\{\begin{matrix}v\\j\end{matrix}\right\}t^j.
\end{align}
If
\[
 d_i(v)=\gcd_{i\leq j\leq v}
          \left((j-1)!\left\{\begin{matrix}v\\j\end{matrix}\right\}\right),
 \qquad d_{v+1}(v)=0,
\]
then $F_\gamma^iK_0=d_i(v)I$ for $1\leq i\leq v$, and $F_\gamma^{v+1}=0$. In particular
\[
 \operatorname{gr}_\gamma^vK_0\simeq\mathbb Z^S,
 \qquad
 \operatorname{gr}_\gamma^{v-1}K_0\simeq(\mathbb Z/2)^S
 \quad\text{when }v\geq3\text{ is odd}.
\]
\end{proposition}
\begin{proof}
Let $B=\mathbb A^v\setminus Y$ and $U=\pi^{-1}(B)$. Localization and homotopy invariance give $\operatorname{CH}^j(B)=0$ for $j>0$. The morphism $U\to B$ factors as the affine-line bundle in $t$ followed by a torsor under the kernel of the surjection $\mathcal O_B^h\xrightarrow{(f_i)}\mathcal O_B$; hence affine-bundle invariance gives $\operatorname{CH}^j(U)=0$ for $j>0$. Localization for the complement $\coprod_{a,e}D_{a,e}$ shows that $\operatorname{CH}^v(X)$ is generated by these components and that all other positive-codimension groups vanish. Flat pullback of the point class $[a]=0$ on $\mathbb A^v$ gives $[D_{a,0}]+[D_{a,1}]=0$.

To show that there are no further integral relations, use Theorem~\ref{thm:cubic} and duality: $H^{2v}(X_{\bar k},\mathbb Q_\lambda)=\mathbb Q_\lambda(-v)^S$. Purity and localization for the smooth closed complement give a surjection from its $2S$ fundamental classes to this group; the target beyond it is $H^{2v}(U)=H^{2v}(B)=0$. The last equality follows from purity for the finite set $Y\subset\mathbb A^v$. The opposite-component relations leave the $S$ classes $u_a$ spanning an $S$-dimensional space, so they are independent. An integral relation would give a rational relation and hence have every coefficient zero. This proves both the integral Chow presentation and the cycle-map assertion. Products of positive classes vanish because their codimension is $2v\neq v$.

For $K$-theory, use the motivic wedge and Bott periodicity over the regular base. Every $v$-fold projective-line smash is a suspension, whose reduced diagonal is null. Thus products of two reduced classes vanish and $I^2=0$. The reduced $K_0$ of each sphere is a copy of $K_0$ of the base, giving $I=\mathbb Z^S$. On $(\mathbb P^1)^v$ write $u_j=[\mathcal O(1)_j]-1$. The projective-bundle formula gives $u_j^2=0$, and the Bott generator of the smash pulls back, up to sign, to $u_1\cdots u_v$. Since Adams operations are ring homomorphisms and act on a line bundle by its $m$th power,
\[
 \psi^m(u_1\cdots u_v)
    =\prod_j((1+u_j)^m-1)=m^v u_1\cdots u_v.
\]
The operation is natural under the motivic equivalence: this is precisely the unstable $\lambda$-ring statement of Riou (2010, Theorem~2.3.1 and Proposition~2.3.2; preprint pp.~10--11) \cite{ff-riou}.

In the torsion-free ring tensored with $\mathbb Q$, Newton's formula reads
\[
 \log\lambda_t(z)=\sum_{m\geq1}\frac{(-1)^{m-1}}m\psi^m(z)t^m.
\]
Every term belongs to the square-zero ideal, so the logarithm is just $\lambda_t(z)-1$. Substitution gives the first formula, whose integral coefficients make it valid already over $\mathbb Z$. Now $\gamma_t(z)=\lambda_{t/(1-t)}(z)$. Its scalar polynomial $P_v$ satisfies
$P_1(t)=t$ and $P_{v+1}(t)=t(1-t)P_v'(t)$. The Stirling recurrence proves the displayed second formula by induction.

By definition the gamma filtration is generated by products of $\gamma^{j}(z)$ with total index at least $i$; two positive-index factors multiply to zero. Therefore it is precisely the gcd of the indicated integer coefficients times $I$, proving the filtration formula. At the top, $d_v=(v-1)!$. For odd $v\geq3$,
\[
 d_{v-1}=\gcd\left((v-2)!\binom v2,(v-1)!\right)
          =\frac{(v-1)!}{2},
\]
so the preceding quotient is $(\mathbb Z/2)^S$. This torsion is in the associated graded, while $K_0$ itself and the ordinary Chow groups are torsion-free.
\end{proof}

For a representation consisting of the verifier and its prescribed cubic construction, the ranks just calculated are $\#\mathrm P$-complete, their nonvanishing is NP-complete, and their vanishing is coNP-complete. For gamma torsion, append forced-zero wire variables until $v$ is odd and at least three; each old accepting word has one extension, so this changes neither $S$ nor the complexity assertion. In particular affine-space recognition and $\mathbb A^1$-contractibility are coNP-complete on this encoding. These membership statements use the recoverable Boolean verifier; they do not assert a general upper bound for arbitrary sparse cubic presentations with unknown provenance.

\subsection{Comparing actual cohomology representations}
\begin{proposition}\label{prop:pp-tate}
Fix a finite field $\mathbb F_q$ and a prime $\lambda$ different from its characteristic. For two cubic schemes constructed from Boolean verifiers over this field and padded to the same $v,h$, deciding whether there exists a Frobenius-equivariant injection
\[
 H_c^{2h}(X_1,\mathbb Q_\lambda)\hookrightarrow
 H_c^{2h}(X_2,\mathbb Q_\lambda)
\]
is PP-complete. Hardness persists when both ranks are positive. This concerns linear maps of cohomology representations, not morphisms of schemes or algebraic correspondences.
\end{proposition}
\begin{proof}
Pad the shorter variable list with forced-zero variables and repeat an existing Boolean equation to equalize the numbers of equations. Neither operation changes its witness count. The two representations are $\mathbb Q_\lambda(-h)^{S_1}$ and $\mathbb Q_\lambda(-h)^{S_2}$, so every linear map commutes with Frobenius and an injection exists exactly when $S_1\leq S_2$.

For hardness, take a verifier on $n\geq1$ primary bits with $a$ accepting words. Construct a right verifier with $S_R=a+1$ by a tagged union with one forced word, and a left verifier with $S_L=2^{n-1}+2$ by a tagged union of a half-cube and two forced distinct words. Then $S_L\leq S_R$ exactly when $a>2^{n-1}$, the PP majority test. Both counts are positive, and all constructions have polynomial size and unique wire extensions.

For membership, let $A,B$ be the two counts after zero-padding to a common witness length $w$. Introduce a bit $r$ and restrict both predicates to $r=0$, so their common domain has size $M=2^{w+1}$ without changing their counts. A selector accepts the second predicate on its first branch and the complement of the first on its second branch. It has $M+B-A$ accepting paths among $2M$. Add a free bit to double the paths, then change one known rejecting path (with second-predicate branch, $r=1$, and all other bits zero) to accepting. The resulting count is $2M+2(B-A)+1$ among $4M$, strictly above half exactly when $A\leq B$. This proves PP membership with equality treated correctly.
\end{proof}

\subsection{A complete two-variable example}
Take $\Phi=(x_1\vee x_2)$, whose satisfying assignments are $(0,1),(1,0),(1,1)$. Its integral Boolean scheme has equations
\[
 f_1=x_1^2-x_1,\qquad f_2=x_2^2-x_2,\qquad
 f_3=(1-x_1)(1-x_2).
\]
No gate variables are necessary, so $v=2$, $h=3$, and $S=3$. The cubic is
\begin{align*}
 t^2-t={}&y_1x_1^2-y_1x_1+y_2x_2^2-y_2x_2\\
         &+y_3-y_3x_1-y_3x_2+y_3x_1x_2.
\end{align*}
Over $\mathbb F_2$, the three satisfying base points each have $2\cdot2^3=16$ points above them, and $(0,0)$ has $2^3=8$, giving $56$. Over $\mathbb F_3$, three base points have $54$ points each and the other six have $27$ each, giving $324$. In general the answer is $q^5+3q^3$. The Frobenius eigenvalues on actual compactly supported \'etale cohomology are $q^3$ with multiplicity three in degree six and $q^5$ with multiplicity one in degree ten. Thus
\[
 Z(X_\Phi/\mathbb F_q,T)=(1-q^3T)^{-3}(1-q^5T)^{-1}.
\]
Its ordinary degree-four cohomology is $\mathbb Q_\lambda(-2)^3$, its Chow ring is square-zero on the three degree-two classes $[D_{a,0}]$, and its complex realization is a wedge of three four-spheres. For $z$ in the reduced $K_0$ group,
$\gamma_t(z)=1+(t-t^2)z$. This single example follows the same three witnesses from the Boolean scheme to the cubic, its points, its Frobenius, and its integral operations.
